\documentclass[../main.tex]{subfiles}
\begin{document}

\begin{center}
    \Large{\textbf{ABSTRACT}}\\
\end{center}
\vspace{0.5cm}

In modern cybersecurity, detecting Living-off-the-Land (LotL) attacks remains a critical challenge due to malicious behaviors blending seamlessly into routine administrative activities. Although unsupervised anomaly detection (UAD) models are highly anticipated for their zero-day threat detection capabilities, their operational performance often degrades severely in practice without clear theoretical explanation. This study addresses this research gap by proposing an independent theoretical framework to quantify the intrinsic difficulty of attack datasets, moving away from empirical, algorithm-dependent evaluations.

This thesis makes three core contributions. First, we propose DensPct-Mahal, an algorithm-independent intrinsic difficulty metric to quantify the camouflage level of LotL attacks. Empirical evidence suggests that this metric explains the performance degradation of anomaly detection algorithms more effectively than the widely used 1D Marginal Overlap (DDO). Second, we construct and systematize a mixed-type feature space, V10, specifically designed for Process Forensics data from Windows Sysmon. Third, we develop a Controlled Injection data generation pipeline simulating the MITRE ATT\&CK lifecycle, producing a rigorous benchmark dataset for the community. Furthermore, the study identifies the existence of Local Pockets and employs a joint non-parametric estimator (KDE\_Joint) as an existence proof to bypass the Global Center Bias. However, we also reveal the critical vulnerabilities of raw KDE on sparse data spaces, thereby directing future research towards more robust neighborhood-based measurement techniques.

\end{document}